\documentclass[%
%handout,%
9pt,hyperref={pdfpagelabels=false},aspectratio=169]{beamer}

\mode<presentation> {
    \usetheme{HHUD}
    \setbeamercovered{invisible}
}


%\listfiles

\usepackage[english]{babel}
\usepackage[utf8x]{inputenc}
\usepackage{times}
\usepackage[T1]{fontenc}
\usepackage{amsmath}
\usepackage{subfigure}
\usepackage{graphicx}
\usepackage{hyperref}
\usepackage{xmpmulti}
\usepackage{multicol}

%\usepackage{appendixnumberbeamer}

% potentially useful packages
\usepackage{nicefrac}
\usepackage{tikz}


% background image
\usebackgroundtemplate{\includegraphics[width=\paperwidth, height=\paperheight]{fig/background_16-9.pdf}}
% commands for low and high decoration in frame foot
\newcommand{\footdecorationlow}{\usebackgroundtemplate{\includegraphics[width=\paperwidth, height=\paperheight]{fig/background_small_16-9}}}
%\newcommand{\footdecorationhigh}{\usebackgroundtemplate{\includegraphics[width=\paperwidth, height=\paperheight]{fig/titel_slide_info.pdf}}}
\newcommand{\footdecorationend}{\usebackgroundtemplate{\includegraphics[width=\paperwidth, height=\paperheight]{fig/end_slide.pdf}}}


% Fix build errors on debian (http://bugs.debian.org/cgi-bin/bugreport.cgi?bug=452333)
\providecommand \thispdfpagelabel[1]{} {}

%% The following lines are used to include section slides at the beginning of a new section;
% \AtBeginSection[] {
%   \footdecorationhigh
%   \begin{frame}<beamer>
%     \thispagestyle{empty}
%     \frametitle{Gliederung}
%     \vspace{-5mm}
%     \tableofcontents[currentsection]
%   \end{frame}
%   \footdecorationlow
% }


% Template based on slides from the Institute of Computer Science
% Modifications regarding the design by Marcel Käufler and Sara Schulte
% % % % % % % % % %  CHANGE TOPIC AND AUTHOR INFORMATION HERE % % % % % % % % %
                    
\title{Entwicklung einer Anwendung zum Vermitteln von Lerninhalten über E-Graphs und Equality Saturation für die Lehre}
%\subtitle{Add further information / subtitles / authors ...}
\author{Ben Steinhauer}
\date{17.02.2025}
\institute{Lehrstuhl für Softwaretechnik und Programmiersprachen \\Heinrich-Heine-Universität Düsseldorf}
\subject{Informatik}

%
% Start of the document
%
\begin{document}
%\footdecorationhigh
\begin{frame}[noframenumbering]
\thispagestyle{empty}
%\titlepage
	%\begin{columns}
	 %  	\column{0.4\paperwidth}
	  % 	{
	   %		\footnotesize
	   	%	\color{hhuUniBlau}
	   		%\put(50,-190){\insertdate}
	   		
	   	%}
	   	%\column{0.6\paperwidth}
	   	\centering 
		\Large\color{black} \textbf{Bachelorarbeit}\\[\baselineskip]
	   	
		\color{hhuUniBlau}
	   	\LARGE \inserttitle\\[\baselineskip]
	   	%\normalsize \insertsubtitle\\[1.2\baselineskip]
	   	
	   	\large \insertauthor\\[\baselineskip]

	   	
	   	\color{black}
	   	\small
	   	\insertinstitute\\[\baselineskip]
	   	\small  \insertdate
	%\end{columns}
\end{frame}

% \footdecorationlow
% \begin{frame}
% 	\thispagestyle{empty}
% 	\frametitle{Gliederung}
% 	Diese Gliederung ist optional und nicht unbedingt empfohlen (vor allem bei Präsentationen unter einer halben Stunde; löschen in master.tex)
% 	\vspace{-5mm}
% 	\tableofcontents
% \end{frame}
% \footdecorationlow

% % % % % % % % % % Start of the conent % % % % % % % % % %

%%% Slides can directly be included here instead of additional tex files
\section{Einleitung}

\subsection{Motivation}

%~\cite{phaseorder-2009}

\subsection{Ziel der Arbeit}

Das Ziel dieser Bachelorarbeit ist es, ein sinnvolles Werkzeug für die Lehre zu erstellen,
um Studentinnen und Studenten die Themen \textbf{E-Graphs} und \textbf{Equality Saturation} näher zu bringen.
Dabei sollen sie die Möglichkeit haben, sich sowohl auf theoretischer als auch praktischer Ebene mit E-Graphs auseinander setzen zu können.
Die theoretische Ebene soll den Studenten die notwendigen Hintergrundkenntnisse vermitteln sowie einen Enblick in die Implementierung geben.
Die praktische Ebene soll Schritt für Schritt aufzeigen, wie der \textbf{E-Graph} aufgebaut wird, und wie an diesem \textbf{Equality Saturation} durchgeführt werden kann.
Für grö{\ss}tmöglichen Nutzen soll die Anwendung plattformunabhängig sein und möglichst nur von \textit{Open-Source-Software} (OSS) Gebrauch machen.
Damit wird das Problem der unterschiedlichen Betriebssysteme der Studenten umgangen und zeitgleich die Hürden für Erweiterungen gesenkt.

\subsection{Aufbau der Arbeit}

\noindent {\itshape Kapitel \ref{sec:grundlagen} - Grundlagen}

Dieses Kapitel zeichnet den Beginn der Arbeit, indem notwendige (mathematische) Kenntnisse über \textbf{E-Graphs} und \textbf{Equality Saturation} vermittelt werden.

\noindent {\itshape Kapitel \ref{sec:entwicklung} - Entwicklung}

\noindent {\itshape Kapitel \ref{sec:entwicklung} - Architektur}

\noindent {\itshape Kapitel \ref{sec:entscheidungenundprobleme} - Entscheidungen und Probleme}

\noindent {\itshape Kapitel \ref{sec:zusammenfassung} - Zusammenfassung}

\subsection{Verwandte Arbeiten}

\noindent\textbf{egg} Das Akronym \textit{egg} steht für \textit{e-graphs good} und ist eine in der Programmiersprache \textit{Rust} geschriebene Bibliothek.
Die Bibliothek implementiert die Datenstruktur E-Graph und stellt Funktionen für die Interaktion mit dieser zur Verfügung.
Neben der Erzeugung von E-Graphs kann auch \textit{Equality Saturation} auf diesen durchgeführt werden.
Das gesamte Projekt wurde in einem Paper vorgestellt, auf welches sich diese Arbeit stützt~\cite{2021-egg}.


\EnableTitleSlide
\section{Ergebnisse \& Diskussion}

\begin{frame}[fragile]{Komplikationen während der Entwicklung}
    \begin{enumerate}
        \item \onslide<1->{\textbf{Kombination zweier Implementierungen}}
        
        \only{
            \vspace{8mm}
            \begin{center}
                \includegraphics[scale=0.75]{utils/colabegg.png}
            \end{center}
        }<1>

        \item \onslide<2->{\textbf{Darstellung der Datenstruktur}}
        
        \only{
            \vspace{2mm}
            \begin{center}
                \includegraphics[scale=0.4]{utils/egraph_exp.png}
            \end{center}
        }<2>

        \begin{onlyenv}<3>
            \vspace{3mm}
            \begin{center}
\begin{minted}[fontsize=\small]{python}    
def egraph_to_dot(self, nodesep=0.5, ranksep=0.5, marked_eclasses = []):
    dot_commands = [
        "digraph parent { graph [compound=true, nodesep=" + str(nodesep)
        + ", ranksep=" + str(ranksep) + "]\n" + """node [fillcolor=white 
        fontname=\"Times-Bold\" fontsize=20 shape=record style=\"rounded, 
        filled\"]\n"""
    ]
    # ... insgesamt 106 Zeilen lang
    return "".join(dot_commands)
\end{minted}
        \end{center}
        \end{onlyenv}

        \item \onslide<4->{\textbf{Spezialfall: Kreis im E-Graph}}
        
        \only<4>{\vspace{3mm}\begin{center}\includegraphics[scale=0.3]{utils/loop.png}\end{center}}
        
\begin{onlyenv}<5>
    \vspace{3mm}
    \begin{center}
\begin{minted}[fontsize=\small]{python}    
def equality_saturation(rules, eterm_id, egraph):
    #
    while True and timeout < 3:
        v = egraph.version
        timeout += 1
        #
        if v == egraph.version:
            break
\end{minted}
\end{center}
\end{onlyenv}

\begin{onlyenv}<6>
    \vspace{3mm}
    \begin{center}
\begin{minted}[fontsize=\small]{python}    
def equality_saturation(rules, eterm_id, egraph):
    #
    while True:
        best_term = _extract_term(eterm_id, egraph)
        if old_term == best_term:
            break
        old_term = best_term
        #  
\end{minted}
\end{center}
\end{onlyenv}
        
        \item \onslide<7->{\color{gray-500}\textbf{Pfadangabe im Server}\color{black}}
        
\begin{onlyenv}<8>
    \vspace{3mm}
    \begin{center}
\begin{minted}[fontsize=\small]{python}    
app.mount("/", 
StaticFiles(directory=realpath(f"{realpath(__file__)}/../static"), html=True),
    name="static"
) # https://github.com/fastapi/fastapi/issues/3550
\end{minted}
\end{center}
\end{onlyenv}

    \end{enumerate}
\end{frame}

\begin{frame}{Ergebnisse}
    \begin{enumerate}
        \item funktionstüchtige Anwendung
        \item plattformunabhängig \& getestet
        \item basiert ausschließlich auf Open-Source-Software
        \item erfüllt Anforderungen des Exposés~\cite{expose}:
        \begin{itemize}
            \color{gray-500}
            \item Benutzeroberfläche im Browser
            \item Erzeugen \& Visualisieren von E-Graphs
            \item Erstellen \& Anwenden von rewrite rules, vordefinierte rewrite rules
            \item Debugging-Feature
            \item Extraktion des optimalen Terms
            \item Export des E-Graphs in gängige Formate
            \item Eingaben als Session abspeichern
            \item Dokumentation
            \color{black}
        \end{itemize}
    \end{enumerate}
\end{frame}

\begin{frame}{Qualität der Software}
    nach ISO/OEC 20510:2011 Standard:\vspace{4mm}

    \begin{enumerate}
        \bf
        \item Übertragbarkeit
        \item Wartbarkeit
        \item Sicherheit
        \item Zuverlässigkeit
        \item Funktionale Eignung
        \item Performance
        \item Kompatibilität
        \item Benutzbarkeit
    \end{enumerate}
\end{frame}

\begin{frame}{Erweiterungen}
    \begin{enumerate}
        \item \textbf{Erweitertes Testing}
        \begin{itemize}
            \item Vergleichsmethode: eigene Implementierung vs. egg 
        \end{itemize}
        \item \textbf{Erstellung von E-Graphs visualisieren}
        \begin{itemize}
            \item schrittweises Darstellung vom Aufbauprozess (AST zu E-Graph) 
        \end{itemize}
        \item \textbf{Abindung an egg}
        \begin{itemize}
            \item Benutzer kann zwischen eigener Implementierung und egg als Backend wählen
            \item benötigt Methoden in Rust
        \end{itemize}
        \item \textbf{E-Class Analysis}
        \begin{itemize}
            \item Conditional and Dynamic Rewrites
            \item Constant Folding
        \end{itemize}
    \end{enumerate}
\end{frame}

\footdecorationlow
\section{Fazit}
\begin{frame}[allowframebreaks]
	\thispagestyle{empty}
	\frametitle{Fazit}
	
\end{frame}


\footdecorationlow
\section{Quellen}
\begin{frame}[allowframebreaks]
	\thispagestyle{empty}
	\frametitle{Literatur}
	Die Quellen werden in der Datei \textit{references.bib} eingefügt und hier angezeigt.
	\nocite{*}
	\bibliography{references.bib}
	\bibliographystyle{abbrv}
\end{frame}

\appendix

\input{anhang}

\end{document}